% !PRINTED 302
\weberpage{302}
Wenn man die erste dieser Formeln $\alpha_1$ mal, die zweite $\alpha_2$ mal
nach einander anwendet u. s. f., so folgt:

\[
\begin{aligned}
\varphi_1(z_1+\alpha_1, z_2, \dots) &\equiv \varphi_1(z_1, z_2, \dots) + \alpha_1\alpha_{1,1} \\
\varphi_1(z_1, z_2+\alpha_2, \dots) &\equiv \varphi_1(z_1, z_2, \dots) + \alpha_2\alpha_{1,2} \\
&\cdots\cdots\cdots\cdots\cdots\cdots\cdots\cdots\cdots,
\end{aligned}
\]

und wenn man in der ersten $z_2$ in $z_2+\alpha_2$ verwandelt, und die
zweite anwendet und so fortfährt, so ergiebt sich schliesslich:

\weq{4}{\begin{aligned}
\varphi_1(z_1+\alpha_1, z_2+\alpha_2, \dots) &\equiv \\
\varphi_1(z_1, z_2, \dots) + \alpha_1\alpha_{1,1} + \alpha_2\alpha_{1,2} + \cdots
\end{aligned}}

Hierin setzen wir nun $z_1, z_2, \cdots = 0$ und schreiben dann an Stelle von
$\alpha_1, \alpha_2, \dots$ wieder $z_1, z_2, \dots$. Ferner bezeichnen
wir $\varphi_1(0, 0, \dots)$ durch $\alpha_1$ und erhalten so aus (4):

\weq{5}{\varphi_1(z_1, z_2, \dots) \equiv \alpha_1 + z_1\alpha_{1,1} + z_2\alpha_{1,2} + \cdots,}

d. h. $\varphi_1$ muss eine \gesperrt{lineare Function} sein.

Genau auf dieselbe Weise verfahren wir mit sämmtlichen Congruenzen (2) und gelangen so zu folgendem Endresultate:

\gesperrt{I. Alle Permutationen der Galois'schen Gruppe $P$ einer primitiven irreduciblen metacyklischen Gleichung vom Grade $p^k$}

\[
\begin{pmatrix} z_1, z_2, \dots, z_k \\ z_1', z_2', \dots, z_k' \end{pmatrix}
\]

\gesperrt{sind von der Form}

\weq{6}{\begin{aligned}
z_1' &\equiv \alpha_{1,1}z_1 + \alpha_{1,2}z_2 + \cdots + \alpha_{1,k}z_k + \alpha_1 \\
z_2' &\equiv \alpha_{2,1}z_1 + \alpha_{2,2}z_2 + \cdots + \alpha_{2,k}z_k + \alpha_2 \quad (\mathrm{mod}\ p). \\
&\cdots\cdots\cdots\cdots\cdots\cdots\cdots\cdots\cdots\cdots
\end{aligned}}

\[
z_k' \equiv \alpha_{k,1}z_1 + \alpha_{k,2}z_2 + \cdots + \alpha_{k,k}z_k + \alpha_k
\]

Das System der Congruenzen (6) stellt immer dann und nur dann eine Permutation dar, wenn die Determinante

\[
\sum \pm \alpha_{1,1}\alpha_{2,2}\cdots\alpha_{k,k}
\]

nicht durch $p$ theilbar ist, weil dann und nur dann auch umgekehrt zu jedem Systeme der $z'$ ein bestimmtes System der $z$ (nach dem Modul $p$) gehört.

Der Inbegriff aller Permutationen (6) ist in der That eine Gruppe, die die \gesperrt{allgemeine lineare Congruenzgruppe} heisst, und in ihr müssen die metacyklischen Gruppen $P$ ent-
%CONT>>
